\section{Tautan \textit{Repository}}

\url{https://github.com/jsndwrd/IF3170_G08_Teazzi}

\section{Pernyataan Penggunaan AI Generatif}

Kelompok menggunakan asisten AI generatif selama pengerjaan tugas ini. Seluruh
keluaran AI telah diperiksa, dijalankan, dan dipahami oleh anggota kelompok. Surat
pernyataan pada subbab berikutnya mengikuti format Panduan Penggunaan
\textit{Generative AI} ITB~\cite{itb2024panduanai}.

\begin{table}[H]
	\centering
	\caption{Penggunaan AI generatif}
	\label{tab:genai}
	\begin{tabular}{p{2.5cm} p{10.7cm}}
		\toprule
		\textbf{Alat} & \textbf{Penggunaan} \\
		\midrule
		Claude & Membantu menyusun kerangka laporan. \\
		\bottomrule
	\end{tabular}
\end{table}

% ---------- Surat pernyataan (format Panduan AI ITB) ----------
% Tanda tangan: simpan gambar di docs/public/ttd/<NIM>.png (latar transparan atau putih).
\newcommand{\ttdanggota}[2]{%
	\begin{minipage}[t]{6.5cm}
		\centering
		\IfFileExists{public/ttd/#2.png}{%
			\includegraphics[height=2cm,width=5.5cm,keepaspectratio]{ttd/#2}%
		}{%
			\rule{0pt}{2cm}%
		}\\[2pt]
		\rule{5.5cm}{0.4pt}\\[2pt]
		#1\\
		#2
	\end{minipage}%
}

\clearpage
\section{Surat Pernyataan Penggunaan AI}
\noindent Dengan ini, kami:

\noindent
\begin{tabular}{@{}l l@{}}
	Nama/NIM            & : 1. Jason Edward Salim/13524034 \\
	                    & \phantom{:} 2. Stefani Angeline Oroh/13524064 \\
	                    & \phantom{:} 3. Bryan Pratama Putra Hendra/13524067 \\
	                    & \phantom{:} 4. Angelina Andra Alanna/13524079 \\
	                    & \phantom{:} 5. Michael James Liman/13524106 \\
	Fakultas/Sekolah    & : Sekolah Teknik Elektro dan Informatika \\
	Kode Mata Kuliah    & : IF3170 \\
	Nama Mata Kuliah    & : Inteligensi Artifisial \\
	Judul Tugas         & : Tugas Besar 2 IF3170 Inteligensi Artifisial 2026/2027 \\
\end{tabular}

\medskip
\noindent menyatakan bahwa kami menggunakan AI dalam pengerjaan maupun penyusunan
luaran tugas pada mata kuliah yang tertulis di atas.

\medskip
\noindent Jika ya, alat AI/kecerdasan buatan yang digunakan adalah:

\begin{center}
\small
\renewcommand{\arraystretch}{1.25}
\begin{tabular}{|>{\raggedright\arraybackslash}p{5.3cm}|>{\centering\arraybackslash}p{2.4cm}|>{\raggedright\arraybackslash}p{5.4cm}|}
	\hline
	\multicolumn{1}{|c|}{\textbf{Lingkup Pekerjaan}} & \textbf{Digunakan?} & \multicolumn{1}{c|}{\textbf{Tingkat Penggunaan}} \\
	\hline
	Pemeriksaan Ejaan: menggunakan \textit{tools} seperti Grammarly, Paperpal,
	DeepL, Quillbot, Grammarbot, LanguageTool, ProWritingAid, ChatGPT, Google
	Gemini, atau sejenisnya.
	& Tidak
	& \\
	\hline
	Pembuatan Teks: menggunakan \textit{tools} seperti ChatGPT, Gemini, Paperpal,
	Copilot, GrammarlyGO, WordAI, WriteSonic, Jasper, Jenni AI, atau sejenisnya.
	& Ya
	& Membantu menyusun kerangka laporan. \\
	\hline
	Bantuan Diskusi dalam Penyusunan Konten: menggunakan \textit{tools} seperti
	ChatGPT, Gemini, Copilot, Perplexity, Jenni AI, Zoom-Companion, atau
	sejenisnya.
	& Ya
	& Mempelajari cara kerja algoritma supervised learning secara mendalam sebelum dikembangkan menjadi barisan kode. \\
	\hline
	Pembuatan Gambar, Video, dan/atau Grafik: menggunakan \textit{tools} seperti
	Craiyon, DALL-E, Midjourney, Stable Diffusion, Microsoft Designer, Gemini,
	Canva AI, atau sejenisnya.
	& Tidak
	& \\
	\hline
	Lainnya, jelaskan:
	& Tidak
	& - \\
	\hline
\end{tabular}
\end{center}

\noindent Kami juga menyatakan bahwa setiap penggunaan AI/kecerdasan buatan yang
dilakukan pada mata kuliah tersebut telah dinyatakan di dalam surat ini.

\bigskip
\begin{flushright}
	Bandung, 23 November 2026
\end{flushright}

\bigskip
\noindent
\begin{center}
	\ttdanggota{Jason Edward Salim}{13524034}\hspace{1cm}%
	\ttdanggota{Stefani Angeline Oroh}{13524064}

	\vspace{0.8cm}
	\ttdanggota{Bryan Pratama Putra Hendra}{13524067}\hspace{1cm}%
	\ttdanggota{Angelina Andra Alanna}{13524079}

	\vspace{0.8cm}
	\ttdanggota{Michael James Liman}{13524106}
\end{center}
